\documentclass[11pt,a4paper]{article}

\usepackage{thga-db}

\renewcommand{\thgaKurs}{Databases}
\renewcommand{\thgaDatum}{Summer Term 2026}

\begin{document}
\thispagestyle{empty}
\thgaTitel{Term Project – User Manual}%
{Wartungsmangement Datenbankanwendung}%

\begin{infobox}
\textbf{Author:} Niklas Stratenwerth  \\
\textbf{Module:} Introduction to Database Management Systems · THGA Bochum \\
\textbf{Stack:} PostgreSQL · FastAPI · Docker Compose · tkinter (\texttt{.deb}) \\
\textbf{Repository:} \url{https://github.com/Niklas0511/DBMS_Project} \\
\end{infobox}

\tableofcontents
\vspace{1em}

% ============================================================
\section{Allgemeines}
% ============================================================

Die Wartungsmanagement Desktopanwendung ist eine Desktopanwendung, die es ermöglicht, Wartungen an Geräten zu planen und durchgeführte 
zu dokumentieren. Sollte nicht bereits ein API-Key in der \texttt{.env} Datei hinterlegt sein, muss dieser vor dem anlegen neuer Wartungen, Geräte, Standorte oder Techniker
über das {\texttt{X-API-Key}} Feld eingegeben werden. Das anzeigen der Daten funktioniert auch ohne diesen API-Key.
Darunter besteht die Anwendung aus sechs Tabs, {\texttt{Statistik}}, {\texttt{Geräte}}, {\texttt{Wartungsplan}},{\texttt{Wartungen}}, {\texttt{Standorte}} und {\texttt{Techniker}}.
Um die Bedienung zu erleichtern wird empfohlen, bei erstmaliger Einrichtung der Anwendung zuerst die Standorte, dann die Techniker und danach die Geräte anzulegen. 
Danach können Wartungen geplant und durchgeführt werden.
Auf einzelnen Tabs wird im Folgenden näher eingegangen.

% ============================================================
\section{Statistik Tab}
% ============================================================

In diesem Tab werden die Statistiken der Wartungen angezeigt. Hierzu muss der Benutzer einen Zeitraum auswählen, dies geschieht
durch die Auswahl von Start- und Enddatum im JJJJ-MM-DD Format. Danach kann über die beiden Knöpfe ausgewählt werden, ob die Statistiken nach Standorten oder nach Technikern
aufgeschlüsselt werden sollen. Bei der Aufschlüsselung nach Standort werden nach ID sortiert die einzelnen Standorte mit der Anzahl derr im Zeitraum durchgeführten Wartungen
und der Anzahl der offenen geplanten Wartungen angezeigt. 
Bei der Aufschlüsselung nach Technikern werden nach ID sortiert die einzelnen Techniker mit der Anzahl der im Zeitraum durchgeführten Wartungen
angezeigt. Die Statistiken werden in einer Tabelle angezeigt.

% ============================================================
\section{Geräte Tab}
% ============================================================
In diesem Tab kann entweder nach spezifischen Geräten über ihre Seriennummer gesucht werden 
oder es können alle Geräte eines Standorts angezeigt werden. Diese Auswahl passiert durch die Textfelder {\texttt{Seriennummer}} 
und {\texttt{Standort-ID}} im Bereich Geräte anzeigen. Nach eintragen der Seriennummer und dem Drücken des Knopfes {\texttt{Gerät suchen}}
 wird das Gerät mit der entsprechenden Seriennummer angezeigt. Wird stattdessen eine Standort-ID eingetragen und der Knopf {\texttt{Geräte am Standort}} 
 gedrückt, werden alle Geräte des Standorts angezeigt.
Unabhängig von dieser Auswahl werden die Informationen zu den Geräteen in der Tabelle darunter angezeigt. 
Hierzu gehören die Seriennummer, der Hersteller, der Gerätetyp, das Wartungsintervall, der Standort und die Nächste geplante Wartung.
Neue Geräte können über die Textfelder im Bereich {\texttt{Neues Gerät}} hinzugefügt werden. 
Hierzu müssen alle Felder ausgefüllt werden und der Knopf {\texttt{Gerät anlegen}} gedrückt werden.

% ============================================================
\section{Wartungsplan Tab}
% ============================================================

In diesem Tab können geplante Wartungen angezeigt und neue Wartungen geplant werden.
Gesteuert, welche geplkanten Wartungen angezeigt werden, wird über den Abschnitt {\texttt{Wartungsplan anzeigen}}.
In diesem Abschnitt kann entweder ein Zeitraum im JJJJ-MM-DD Format ausgewählt werden, oder es kann eine Techniker-ID eingetragen werden.
Um dann den Wartungsplan anzuzeigen, muss für den Zeitraum der Knopf {\texttt{Zeitraum laden}} gedrückt werden und 
für die Techniker-ID der Knopf {\texttt{Techniker laden}} es ist außerdem möglich über den Knopf {\texttt{überfällige Laden}} alle
Wartungen anzuzeigen, die noch nicht durchgeführt wurden und deren geplantes Datum überschritten wurde.

Unabhängig von der Auswahl werden die geplanten Wartungen in der Tabelle darunter angezeigt. Hierzu gehören die ID der Wartung, das geplante Datum, die Seriennummer des Geräts,
die Techniker-ID, der Name des Technikers, der Standort und der Status.
Neue Wartungen können über den Abschnitt {\texttt{Wartung planen oder ändern}} geplant werden.
Um einen neue Wartung zu planen, kann das Feld {\texttt{Plan-ID}} leer gelassen werden, es muss nur das geplante Datum, die Seriennummer des Geräts und die Techniker-ID eingetragen werden.
Um eine bestehende Wartung zu ändern, muss die ID der Wartung in das Feld {\texttt{Plan-ID}} eingetragen werden und die Felder, die geändert werden sollen, müssen ausgefüllt werden.
Danach kann entweder der Knopf {\texttt{Planen}}, der Knopf {\texttt{Ändern}} oder der Knopf {\texttt{Löschen}} gedrückt werden, je nachdem ob eine neue Wartung geplant oder eine bestehende 
Wartung geändert oder gelöscht werden soll. 

% ============================================================
\section{Wartungen Tab}
% ============================================================

In diesem Tab können durchgeführte Wartungen angezeigt und neue Wartungen dokumentiert werden.
Gesteuert, welche durchgeführten Wartungen angezeigt werden, wird über den Abschnitt {\texttt{Wartungshistorei}}.
Hier kann im Feld {\texttt{Seriennummer}} die Seriennummer des Geräts eingetragen werden, dessen Wartungshistorie angezeigt werden soll.
Danach muss der Knopf {\texttt{Historie laden}} gedrückt werden, um die Wartungshistorie anzuzeigen. Die Wartungshistorie wird in 
der Tabelle darunter angezeigt. Hierzu gehören die ID der Wartung, das Datum der Wartung, die Seriennummer des Geräts, die Techniker-ID, der
Name des Technikers, optional die Seriennummer eines Ersatzgeräts und die zugehörige Plan ID.
Neue Wartungen können über den Abschnitt {\texttt{Durchgeführte Wartung dokumentieren}} dokumentiert werden. 
Hierzu müssen das Datum der Wartung, die Seriennummer des Geräts, die Techniker-ID, optional die Seriennummer eines Ersatzgeräts und optional die zugehörige Plan ID eingetragen werden.
Danach muss der Knopf {\texttt{Wartung speichern}} gedrückt werden, um die Wartung zu speichern.

% ============================================================
\section{Standorte Tab}
% ============================================================

In diesem Tab können Standorte angezeigt und neue Standorte angelegt werden.
Alle Standorte werden in der Tablle {\texttt{Vorhandene Standorte}} angezeigt. Hierzu gehören die ID des Standorts,
die Postleitzahl und die Adresse. Diese Liste kann bei Bedarf über den Knopf {\texttt{Aktualisieren}} aktualisiert werden. Außerdem kann nach
anklicken eines Standortes in der Tabelle, über den Knopf {\texttt{Auswahl für Gerät übernehmen}} die ID des Standorts in das Feld {\texttt{Standort-ID}} im Abschnitt {\texttt{Geräte}}
übernommen werden, um ein neues Gerät an diesem Standort anzulegen. Neue Standorte können über den Abschnitt {\texttt{Neuen Standort anlegen}} angelegt werden.
Hierfür müssen die Postleitzahl und die Adresse eingetragen werden und der Knopf {\texttt{Standort anlegen}} gedrückt werden.

% ============================================================
\section{Techniker Tab}
% ============================================================

In diesem Tab können Techniker angezeigt und neue Techniker angelegt werden.
Alle Techniker werden in der Tablle {\texttt{Vorhandene Techniker}} angezeigt. Hierzu gehören die ID des Technikers, der Name, die Standort IDs und die Gerätetypen.
Diese Liste kann bei Bedarf über den Knopf {\texttt{Aktualisieren}} aktualisiert werden. Außerdem kann nach anklicken eines Technikers in der Tabelle,
über den Knopf {\texttt{ID in beide Formulare übernehmen }} die ID des Technikers in das Feld {\texttt{Techniker-ID}} im Abschnitt {\texttt{Wartungsplan}} und 
in das Feld {\texttt{Techniker-ID}} im Abschnitt {\texttt{Wartungen}} eingetragen werden, um eine neue Wartung zu planen oder eine durchgeführte Wartung zu dokumentieren.
Neue Techniker können über den Abschnitt {\texttt{Techniker und Zuständigkeiten}} angelegt werden.
Hierfür muss der Name des Technikers in das Feld {\texttt{Name}} eingetragen werden. 
Darunter können unter {\texttt{Zuständige Standorte}} alle Standorte ausgewählt werden, für welche der Techniker zuständig ist.
Außerdem können unter {\texttt{Gerätetypen}} alle Gerätetypen mit Semikolon getrennt eingetragen werden, für welche der Techniker zuständig ist.
Darunter kann der Knopf {\texttt{Speichern}} gedrückt werden, um den Techniker anzulegen. oder der Knopf {\texttt{Neu/Eingabe leeren}} gedrückt werden,
um die Eingabefelder zu leeren.

\end{document}(example)